% !TeX program = xelatex
\documentclass[11pt,en]{elegantpaper}
\usepackage{amsmath, amssymb, amsthm}
\title{MATH4665 Homework 3}
\author{Jiahan XIA}

\begin{document}
\maketitle 

\section*{Question 1}
\textbf{Show that for any measure space $(X, \mathcal{A}, \mu)$ and $1 \le p \le \infty$, the normed space $(L^p(X), \|\cdot\|_p)$ is complete.}

\begin{proof}
We divide the proof into two cases: $1 \le p < \infty$ and $p = \infty$.

\textbf{Case 1: $1 \le p < \infty$} \\
Let $\{f_n\}$ be a Cauchy sequence in $L^p(X)$. Our goal is to find a function $f \in L^p(X)$ such that $\|f_n - f\|_p \to 0$ as $n \to \infty$. 

Because $\{f_n\}$ is Cauchy, we can extract a rapidly Cauchy subsequence $\{f_{n_k}\}$ such that for all $k \ge 1$:
$$\|f_{n_{k+1}} - f_{n_k}\|_p < \frac{1}{2^k}$$
Define the functions $g_K(x) = \sum_{k=1}^K |f_{n_{k+1}}(x) - f_{n_k}(x)|$ and $g(x) = \sum_{k=1}^\infty |f_{n_{k+1}}(x) - f_{n_k}(x)|$. 
By Minkowski's Inequality, for any $K$:
$$\|g_K\|_p \le \sum_{k=1}^K \|f_{n_{k+1}} - f_{n_k}\|_p < \sum_{k=1}^K \frac{1}{2^k} < 1$$
By the Monotone Convergence Theorem (or Fatou's Lemma), since $g_K \uparrow g$ almost everywhere (a.e.), we have:
$$\|g\|_p \le 1 < \infty$$
This implies that $g \in L^p(X)$, and therefore $g(x) < \infty$ almost everywhere. Thus, the series 
$$f_{n_1}(x) + \sum_{k=1}^\infty (f_{n_{k+1}}(x) - f_{n_k}(x))$$
converges absolutely for almost every $x \in X$. Let $f(x)$ be its limit. Note that the $K$-th partial sum of this series telescopes to $f_{n_{K+1}}(x)$, meaning $f_{n_k}(x) \to f(x)$ pointwise a.e. as $k \to \infty$.

We now show that $f \in L^p(X)$ and $f_n \to f$ in $L^p$. Since $\{f_n\}$ is Cauchy, for any $\epsilon > 0$, there exists an $N$ such that for all $n, m \ge N$, $\|f_n - f_m\|_p < \epsilon$. 
Using Fatou's Lemma on the sequence $k \mapsto |f_{n_k} - f_m|^p$ (fixing $m \ge N$), we get:
$$\int_X |f - f_m|^p \, d\mu \le \liminf_{k \to \infty} \int_X |f_{n_k} - f_m|^p \, d\mu \le \epsilon^p$$
This shows that $f - f_m \in L^p(X)$, and since $f_m \in L^p(X)$, $f = (f - f_m) + f_m \in L^p(X)$. Furthermore, $\|f - f_m\|_p \le \epsilon$ for all $m \ge N$, which means the original sequence $f_m$ converges to $f$ in $L^p$.

\textbf{Case 2: $p = \infty$} \\
Let $\{f_n\}$ be a Cauchy sequence in $L^\infty(X)$. For every $n, m \in \mathbb{N}$, there exist null sets $A_n$ and $B_{n,m}$ such that $|f_n(x)| \le \|f_n\|_\infty$ on $A_n^c$ and $|f_n(x) - f_m(x)| \le \|f_n - f_m\|_\infty$ on $B_{n,m}^c$. Let $E = \left(\bigcup_n A_n\right) \cup \left(\bigcup_{n,m} B_{n,m}\right)$. The countable union of null sets is a null set, so $\mu(E) = 0$.

On the set $X \setminus E$, the sequence $\{f_n\}$ is uniformly Cauchy. Therefore, it converges uniformly to some bounded function $f$. We define $f(x) = 0$ for $x \in E$. 
Because the convergence is uniform on $X \setminus E$, we have $\|f_n - f\|_\infty \to 0$ as $n \to \infty$. Thus, $L^\infty(X)$ is complete.
\end{proof}

\vspace{0.5cm}

\section*{Question 2}
\textbf{Prove that if $p \neq 2$, no inner product can generate the usual $L^p$ norm.}

\begin{proof}
A norm $\|\cdot\|$ on a vector space is generated by an inner product if and only if it satisfies the Parallelogram Law:
$$\|f + g\|^2 + \|f - g\|^2 = 2(\|f\|^2 + \|g\|^2)$$
for all $f, g$ in the space. To show that the $L^p$ norm does not come from an inner product when $p \neq 2$, we must find two functions in $L^p(X)$ that violate this identity.

Assume the measure space $(X, \mathcal{A}, \mu)$ contains at least two disjoint measurable sets $A$ and $B$ with finite, strictly positive measure (if the space consists of only a single atom, then $L^p$ is 1-dimensional and trivially an inner product space for all $p$, so we assume the space is not trivial).

Let $f = c_1 \chi_A$ and $g = c_2 \chi_B$, where $c_1, c_2 > 0$ are chosen such that $\|f\|_p > 0$ and $\|g\|_p > 0$. Because $A$ and $B$ are disjoint, we have $f(x)g(x) = 0$ for all $x \in X$.
Consequently:
$$|f(x) + g(x)|^p = |f(x)|^p + |g(x)|^p$$
$$|f(x) - g(x)|^p = |f(x)|^p + |g(x)|^p$$

Evaluating the $L^p$ norms, we get:
$$\|f \pm g\|_p^p = \int_X (|f|^p + |g|^p) \, d\mu = \|f\|_p^p + \|g\|_p^p$$
Thus, $\|f + g\|_p = \|f - g\|_p = \left( \|f\|_p^p + \|g\|_p^p \right)^{1/p}$.

Now, substitute these into the Parallelogram Law. Let $u = \|f\|_p^p$ and $v = \|g\|_p^p$. The left side of the Parallelogram Law becomes:
$$\|f + g\|_p^2 + \|f - g\|_p^2 = (u + v)^{2/p} + (u + v)^{2/p} = 2(u + v)^{2/p}$$
The right side of the Parallelogram Law becomes:
$$2(\|f\|_p^2 + \|g\|_p^2) = 2\left( u^{2/p} + v^{2/p} \right)$$

For the Parallelogram Law to hold, we must have:
$$2(u + v)^{2/p} = 2\left( u^{2/p} + v^{2/p} \right)$$
$$(u + v)^{2/p} = u^{2/p} + v^{2/p}$$
This algebraic identity must hold for all $u, v > 0$. However, it is a well-known fact that $(u + v)^\alpha = u^\alpha + v^\alpha$ for strictly positive $u, v$ if and only if $\alpha = 1$. 

In our case, $\alpha = \frac{2}{p}$. Therefore, the equality holds if and only if $\frac{2}{p} = 1$, which implies $p = 2$. 

If $p \neq 2$, the Parallelogram Law fails for these specifically constructed disjoint functions $f$ and $g$. Hence, the $L^p$ norm cannot be generated by an inner product when $p \neq 2$.
\end{proof}

\end{document}
